\chapter{Used technologies}
\label{technologies}

This chapter briefly introduces the~most important technologies used during
the~development of modules for GRASS GIS. It means that besides GRASS GIS itself, 
the~Python language and its libraries TensorFlow and Keras will be introduced.

\section{GRASS GIS}
\label{grass}

\begin{figure}[H]
   \centering
	\includegraphics[width=0.4\linewidth]{./pictures/grass-logo.png}
	\caption[GRASS GIS logo]{GRASS GIS logo, source: \url{https://grass.osgeo.org/download/logos/}}
      \label{fig:grass-logo}
\end{figure}

The history of \zk{GRASS} \zk{GIS} (an acronym for Geographic Resources Analysis 
Support System Geographic Information System) started in the~year 1982 at the~U.S. 
Army Corps of Engineers Construction Engineering Research Laboratory, but the~
version 1.0 under the~name \zk{GRASS} was released later, in 1985. However, its 
way lead from the~national project to the~international Open GRASS Consortium in 
the mid-1990s. This consortium is perceived as an ancestor of today's Open 
Geospatial Consortium (\zk{OGC}). In 1999, \zk{GRASS} \zk{GIS} was published 
under GNU General Public License (\zk{GPL}).

\zk{GRASS} \zk{GIS} has grown to a~cross-platform free and open-source \zk{GIS} 
maintained by an international team of developers and users and licensed under 
the GNU \zk{GPL} license allowing users to perform geospatial data management 
and analysis for both raster and vector data, image processing, geocoding and
%% ML: v soucasnosti pres 400
%% OP: Dik, cerpal jsem jen z uvodnich informaci GRASSu namisto vlastnich poctu
visualization. Within its more than 400 modules, it supports also 
spatio-temporal data.

%% ML: original -> official ?
For more information about \zk{GRASS} \zk{GIS}, please see its official 
website\footnote{\url{https://grass.osgeo.org/documentation/general-overview/}}.

\section{Python}
\label{python}

\begin{figure}[H]
   \centering
	\includegraphics[width=0.8\linewidth]{./pictures/python-logo.png}
	\caption[Python logo]{Python logo, source: \url{https://www.python.org/community/logos/}}
      \label{fig:python-logo}
\end{figure}

Python is a~cross-platform open-source interpreted programming language. Python 
was invented by a~Dutch programmer and Monty Python's Flying Circus fan Guido 
van~Rossum in the~year 1991. Due to its user-friendly and readable structure, it is 
more and more popular.

It is also one of the most popular languages in the~field of (nay) open-source
\zk{GIS}. Many modules for \zk{GRASS} \zk{GIS} are written in Python and Python
is the~most used language
%% ML: QGIS plugins ML: nejen, GDAL na Python API, Fiona, RasterIO ML:
%% skoro bych to prepsal v tom smyslu, ze vetsina opensource GIS
%% projektu Python podporuje, a nejen open source, napr. pro ESRI je
%% to hlavni skriptovaci platforma na desktopu
in QGIS\footnote{\url{https://qgis.org/}} plugins. Also libraries like
Fiona\footnote{\url{http://toblerity.org/fiona/}} or
RasterIO\footnote{\url{http://rasterio.readthedocs.io/}} are Python-based, and
GDAL\footnote{\url{http://www.gdal.org/}} has its Python API.

For information about programming in Python, please see \cite{diveintopython}.

\section{TensorFlow}
\label{tf}

\begin{figure}[h]
   \centering
	\includegraphics[scale=0.5]{./pictures/tf-logo.png}
	\caption[TensorFlow logo]{TensorFlow logo, source: \url{https://www.tensorflow.org/images/tf_logo_transp.png}}
      \label{fig:tf-logo}
\end{figure}

TensorFlow is an open-source software library developed by Google Brain 
Team\footnote{\url{https://research.google.com/teams/brain/}} primarily for
the~purpose of neural network research. TensorFlow uses data flow graphs, where
a~node represents a~mathematical operation and an edge represents a~tensor
(a~multidimensional array). A~toolkit for visualization of these graphs is called 
TensorBoard.

TensorFlow is available for Python and C++ programming languages. Its usage 
optimizes mathematical expressions and as it was developed for the~purpose of 
neural networks, the~main focus was dedicated to this field and TensorFlow 
provides a~lot of functions useful for deep learning.

A good starting point is the~book Getting Started with TensorFlow \cite{tf}.
%% ML: official
A~documentation and examples can be found also at the~official 
website\footnote{\url{https://www.tensorflow.org/}}.

\section{Keras}
\label{keras}

\begin{figure}[H]
   \centering
	\includegraphics[width=0.8\linewidth]{./pictures/keras-logo.png}
	\caption[Keras logo]{Keras logo, source: \url{https://keras.io/}}
      \label{fig:keras-logo}
\end{figure}

Keras is an open-source software library written in Python and using TensorFlow, 
Microsoft Cognitive Toolkit, Theano or MXNet as a~backend tensor manipulation 
library.

Keras is again developed for the~purpose of deep learning. Its powerful feature 
is the~content of extensible objects for defining different types of layers (see 
chapter \ref{layers}), models, loss functions and other tools widely used in
the~field of \zk{ANN}s.

For a~wider documentation, please see the~official 
website\footnote{\url{https://keras.io/}}.
